\documentclass{article}

% Pre-print submission style for NeurIPS 
\usepackage[preprint,nonatbib]{neurips_2023} % Using standard style

\usepackage[utf8]{inputenc} % allow utf-8 input
\usepackage[T1]{fontenc}    % use 8-bit T1 fonts
\usepackage{hyperref}       % hyperlinks
\usepackage{url}            % simple URL typesetting
\usepackage{booktabs}       % professional-quality tables
\usepackage{amsfonts}       % blackboard math symbols
\usepackage{amsmath}        % basic math
\usepackage{nicefrac}       % compact symbols for 1/2, etc.
\usepackage{microtype}      % microtypography
\usepackage{xcolor}         % colors
\usepackage{graphicx}       % images

\title{The Viability Principle: Information Conservation and Optimal Incoherence in Resource-Constrained Intelligence}

\author{%
  Dr. Rony Charlier (MDL Lab) \\
  Independent Researcher \\
  \texttt{rony.charlier@mdlynor.ai} \\
}

\begin{document}

\maketitle

\begin{abstract}
The dominant paradigm in contemporary deep learning equates intelligence with the unbounded minimization of predictive error (e.g., maximizing log-likelihood, minimizing Kullback-Leibler divergence). However, under finite energetic and structural computational constraints, the pursuit of zero predictive error may lead to exponentially scaling algorithmic costs. In this paper, we introduce the \textbf{MDL Ynor Framework}, proposing that intelligence can be modeled as a variational process that maximizes \textit{Information Viability}. By defining the Ynor Action ($\mathcal{S}_{\text{Ynor}}$) as an integral functional that jointly optimizes effective task-information against both logical divergence and physical computation limits, we provide a unified objective for adaptive systems. Crucially, we formulate the \textit{Optimal Incoherence Theorem}: for any bounded system, maintaining local incoherence ($D_{KL} > 0$) emerges not as an algorithmic flaw, but as a structural necessity for survival under constraint. Empirical validations demonstrate that endogenous computation-scaling governed by the Ynor functional strictly dominates standard fixed-compute baselines on the Pareto frontier.
\end{abstract}

\section{Introduction}

The foundational architecture of contemporary large-scale artificial intelligence models, notably Transformer-based Large Language Models (LLMs), relies heavily on a static inference paradigm. During the forward pass, these systems expend a fixed quantum of computational resources irrespective of the intrinsic difficulty of the input sequence. Current large-scale models do not explicitly optimize computational cost during inference, which may lead to suboptimal behavior under resource constraints.

Current theoretical frameworks addressing AI dynamics generally fall into two camps. The Free Energy Principle (FEP) \cite{friston2010} postulates that adaptive agents act to minimize variational free energy. Conversely, the Information Bottleneck (IB) method \cite{tishby1999} formalizes the extraction of relevant signal by compressing the input space. Yet, neither framework explicitly treats the kinetic physical cost of computation as a continuous endogenous variable subject to active optimization.

To bridge this gap, we introduce the \textbf{MDL Ynor Framework}. We propose a paradigm shift where the metric of an intelligence's generality is uncoupled from raw representational capacity. Instead, we propose that this functional provides a principled objective for resource-constrained intelligence, defining adaptive systems by their capacity to extr\'emalise the Ynor Action ($\mathcal{S}_{\text{Ynor}}$).

\section{Formalism: Ynor State Representation and Variational Action}

\subsection{System Model}
We formalize an adaptive intelligent system as a controlled stochastic process. Let $\Sigma = \langle \mathcal{X}, \mathcal{Y}, \Theta, \Pi, \mathfrak{E}, \mathfrak{T} \rangle$, where $\mathcal{X}, \mathcal{Y}$ denote input and output spaces, $\Theta \subset \mathbb{R}^d$ is the parameter manifold, $\Pi$ is the space of control policies $\pi_t(u_t \mid x_t, \theta_t)$, and $\mathfrak{E}, \mathfrak{T}$ denote environment and task classes. The system evolves according to the stochastic differential equation:
\begin{equation}
d\theta_t = f(\theta_t, \pi_t, \xi_t) dt + \sigma(\theta_t) dW_t
\end{equation}

\subsection{Action Functional and Information Components}
We define three fundamental quantities:
1. \textbf{Effective Information:} $I_{eff}(t) = I(Y_t; X_t) - I(Y_t; X_t \mid T)$, measuring task-relevant extraction.
2. \textbf{Divergence:} $D_{KL}(t) = D_{KL}(P_{\theta_t}(\cdot \mid x_t) \parallel P_t^\star(\cdot \mid x_t))$, the inference error relative to the optimal reference posterior.
3. \textbf{Computational Cost:} $C(t) = C(\pi_t, \theta_t)$, treated as a measurable resource functional.

The Ynor Action is defined as a trajectory-dependent functional:
\begin{equation}
\mathcal{S}_{\text{Ynor}}^\pi(\theta_0; E, T) = \mathbb{E}^\pi \left[ \int_{t_0}^{t_1} \left( I_{eff}(t) - \alpha D_{KL}(t) - \beta C(t) \right) dt \right]
\end{equation}

\subsection{Theorem 1: Necessary Incoherence under Constraints}
\textit{Assume irreducible noise $\mathrm{Var}(\xi_t) > 0$, bounded cost $C(t) \le C_{max}$, and non-zero cost penalty $\beta > 0$. Any admissible policy $\pi$ satisfying viability ($\mathcal{S}_{\text{Ynor}}^\pi \ge \kappa$) must maintain a strictly positive divergence lower bound: $\inf_t D_{KL}(P_{\theta_t} \parallel P_t^\star) \ge \varepsilon$ for some $\varepsilon > 0$.}

\textbf{Proof Sketch:} Reducing $D_{KL} \to 0$ requires arbitrarily precise estimation under stochastic noise. By information-theoretic bounds, this implies increasing model capacity, leading to super-linear growth in computational cost. Since $C(t)$ is bounded, minimizing divergence beyond a threshold violates the viability bound.

\section{Empirical Validation}
We deploy a simulated recurrent inference environment subject to exogenous noise. Monte Carlo evaluations ($N=100$) confirm that the Ynor-controlled endogenous policy actively scales computation based on task difficulty. 
For trivial tasks, the system halts computation early, saving $86\%$ computation vs static baselines. Under an extreme Resource Starvation Test (increasing $\beta$), static baselines suffer mathematical systemic collapse, while the Ynor agent curtails its compute, deliberately accepting a predefined spike in inference error to remain viable ($\mu > 0$).

\section{Discussion and Conclusion}
We propose a shift in the definition of intelligent systems: intelligence optimally maintains information viability under constraint. By introducing a unified variational framework integrating information extraction, approximation error, and computational cost, we provide a principled objective for adaptive systems.
While preliminary, these results suggest that incorporating endogenous computation into objective functions may be a key step toward more general and efficient intelligent systems.

\end{document}
